\section{Theorems}
\hrule \medskip

\begin{theorem}[\textbf{The Modulus and Argument of Complex Numbers}]
    Let $z = a + bi$ be a complex number with $a = Re(z)$ and $b = Im(z)$. That means $a$ representing the real axis and $b$ representing the imaginary axis. Let $(r, \theta)$ be a polar representation of the point with rectangular coordinates $(a, b)$ where $r \geq 0$.
    
    \begin{itemize}
        \item The \textbf{modulus} of $z$, denoted $|z|$, is defined by $|z| = r$.
        \item The angle $\theta$ is an \textbf{argument} of $z$. The set of all arguments of $z$ is denoted $arg(z)$.
        \item If $z \neq 0$ and $-\pi < \theta \leq \pi$, then $\theta$ is the \textbf{principal argument} of $z$, written $\theta = Arg(z)$.
    \end{itemize}
\end{theorem}

\medskip \hrule \medskip 

\begin{theorem}[\textbf{Properties of the Modulus}]
    Let $z$ and $w$ be complex numbers.
    
    \begin{itemize}
        \item $|z|$ is the distance from $z$ to $0$ in the complex plain. 
        \item $|z| \geq 0$ and $|z| = 0$ if and only if $z = 0$
        \item $|z| = \sqrt{Re(z)^2 + Im(z)^2}$
        \item \textbf{Product Rule:} $|zw| = |z||w|$
        \item \textbf{Power Rule:} $|z^n| = |z|^n$ for all natural numbers, $n$
        \item \textbf{Quotient Rule:} $\lvert \frac{z}{w} \rvert = \frac{|z|}{|w|}$, provided $w \neq 0$
    \end{itemize}
\end{theorem}

\medskip \hrule \medskip

\begin{theorem}[\textbf{Properties of the Argument}]
    Let $z$ be a complex number. 
    
    \begin{itemize}
        \item If $Re(z) \neq 0$ and $\theta \in arg(z)$, then $\tan(\theta) = \frac{Im(z)}{Re(z)}$.
        \item If $Re(z) = 0$ and $Im(z) > 0$, then $arg(z) = \{ \frac{\pi}{2} + 2\pi k | k$ is an integer $\}$.
        \item If $Re(z) = 0$ and $Im(z) < 0$, then $arg(z) = \{ -\frac{\pi}{2} + 2\pi k | k$ is an integer $\}$.
        \item If $Re(z) = Im(z) = 0$, then $z = 0$ and $arg(z) = (- \infty, \infty)$.
    \end{itemize}
\end{theorem}

\medskip \hrule \medskip

\begin{theorem}[\textbf{A Polar Form of a Complex Number}]
    Suppose $z$ is a complex number and $\theta \in arg(z)$. The following expression is called a polar form for $z$:
    \[ |z|cis(\theta) = |z| [\cos(\theta) + i \sin(\theta)] \]
\end{theorem}

\medskip \hrule 